\section*{第 \ref{ch:g7:oxygen-in-the-water} 章——水中的氧气}

\begin{solution}{exo:g7:oxygen-in-the-water:1}
每升水中约有 \qty{10}{mg} \omterm{def:g7:what-is-respiration:respiration}{氧气}——大约只有
一升空气所载\omterm{def:g7:what-is-respiration:respiration}{氧气}的三十分之一。
\end{solution}

\begin{solution}{exo:g7:oxygen-in-the-water:2}
上方空气在水面溶入水中，
在水被搅动、拍打、起泡之处最快；
下方的\omterm{def:g5:food-webs:roles}{生产者}白天释放\omterm{def:g7:what-is-respiration:respiration}{氧气}，
在光照充足、绿色繁茂之处最快。
\end{solution}

\begin{solution}{exo:g7:oxygen-in-the-water:3}
\omterm{def:g1:animals-around-us:animal}{动物}的\omterm{def:g7:what-is-respiration:respiration}{呼吸作用}，\omterm{def:g1:plants-around-us:plant}{植物}在夜间的\omterm{def:g7:what-is-respiration:respiration}{呼吸作用}，
以及\omterm{def:g6:decomposers-and-soil:decomposers}{分解者}的\omterm{def:g7:what-is-respiration:respiration}{呼吸作用}——死亡物质丰富时，
其耗氧量会急剧增大。
\end{solution}

\begin{solution}{exo:g7:oxygen-in-the-water:4}
水温升高时，水所能\omterm{def:g7:oxygen-in-the-water:dissolved}{溶解}的\omterm{def:g7:what-is-respiration:respiration}{氧气}减少（溶氧上限下降），
而所有\omterm{def:g7:what-is-respiration:respiration}{呼吸}者的代谢加快（需求
上升）：供给减少，需求增加，两者同时出现。
\end{solution}

\begin{solution}{exo:g7:oxygen-in-the-water:5}
白天，\omterm{def:g5:food-webs:roles}{生产者}的光合作用产氧速度快于
\omterm{def:g7:what-is-respiration:respiration}{呼吸}消耗的速度——曲线一路上升，至傍晚前
较晚时段达到峰值。夜间，所有\omterm{def:g7:what-is-respiration:respiration}{呼吸}者都在
消耗库存\omterm{def:g7:what-is-respiration:respiration}{氧气}，不再有任何产氧——曲线滑落至
日出前的最低点。
\end{solution}

\begin{solution}{exo:g7:oxygen-in-the-water:6}
富氧水域：石块下有鳟鱼（或鳟鱼苗）和
蜉蝣幼虫。贫氧水域：仅有耐受性强的\omterm{def:g5:biodiversity:species}{物种}——
鲤鱼、池塘螺，以及带\omterm{def:g7:what-is-respiration:respiration}{呼吸}管、从水面取气的幼虫。
\end{solution}

\begin{solution}{exo:g7:oxygen-in-the-water:7}
夜间的\omterm{def:g7:what-is-respiration:respiration}{呼吸作用}——\omterm{prop:g3:sorting-living-things:groups}{鱼类}、\omterm{def:g1:plants-around-us:plant}{植物}和淤泥中的
\omterm{def:g6:decomposers-and-soil:decomposers}{分解者}共同消耗——使池塘\omterm{def:g7:what-is-respiration:respiration}{氧气}降至最低，
就在日出前；因此负荷过重的池塘会
在黎明时使\omterm{prop:g3:sorting-living-things:groups}{鱼类}窒息。补救措施：从供氧来源1入手，搅动水面
——整夜开动喷泉或跌水；从供氧来源2入手，
清理过度生长的水生\omterm{def:g1:plants-around-us:plant}{植物}和淤泥，使夜间
\omterm{def:g7:what-is-respiration:respiration}{呼吸作用}减弱（并让白天的产氧作用惠及
整个水体，而不只是水面的浮膜）。
\end{solution}

\begin{solution}{exo:g7:oxygen-in-the-water:8}
废物输入——\omterm{def:g6:decomposers-and-soil:decomposers}{分解者}大量进食并迅速繁殖——
它们集中的\omterm{def:g7:what-is-respiration:respiration}{呼吸作用}消耗\omterm{def:g7:what-is-respiration:respiration}{氧气}——
要求高的\omterm{def:g5:biodiversity:species}{物种}消失，耐受性强的\omterm{def:g5:biodiversity:species}{物种}取代它们。
下游的废物最终被分解，浅滩使水体重新充氧，
要求高的\omterm{def:g5:biodiversity:species}{物种}重新出现：消耗之"漏洞"已堵上，
供氧来源也已跟上。
\end{solution}

\begin{solution}{exo:g7:oxygen-in-the-water:9}
探针只记录一个瞬间；居民记录的是
历史。要求高的\omterm{def:g5:biodiversity:species}{物种}只有在水质
整个停留期间都保持良好时，才会在该站点生存——
包括黎明时分与糟糕的几周，这些是
任何单次采样都捕捉不到的。普查记录的是
一季中的最低值，写在\omterm{def:g1:animals-around-us:animal}{动物}身上。
\end{solution}

\begin{solution}{exo:g7:oxygen-in-the-water:10}
豌豆汤般的绿色是由大量\omterm{def:g6:cell-unit-of-life:unicellular}{单细胞}绿藻
组成的种群：白天它们产氧，但每个温暖的夜晚
整锅“汤”都在\omterm{def:g7:what-is-respiration:respiration}{呼吸}——炎热无风的一周使溶氧上限降低，
并阻止水面补充\omterm{def:g7:what-is-respiration:respiration}{氧气}。黎明时的最低值正走向
零：这是\omterm{prop:g3:sorting-living-things:groups}{鱼类}死亡的天气。整夜及凌晨时段
尽量全力开足喷泉——白天的产氧问题
自然会自行解决。
\end{solution}

\begin{solution}{exo:g7:oxygen-in-the-water:11}
鳟鱼是要求高的\omterm{def:g5:biodiversity:species}{物种}：寒冷的河流能提供其所需的
高溶氧上限；堰坝和激起水花的水道使水面
持续补氧，以抵消密集鱼群的消耗。一个
温暖、静止的湖泊在溶氧上限和补氧两方面同时不足——
那是鲤鱼的地盘，而不是鳟鱼的。
\end{solution}

\begin{solution}{exo:g7:oxygen-in-the-water:12}
货币：每种水下\omterm{def:g7:what-is-respiration:respiration}{呼吸}者都动用同一份
\omterm{def:g7:oxygen-in-the-water:dissolved}{溶解氧}存量——何处能够兴旺，以\omterm{def:g7:what-is-respiration:respiration}{氧气}计价
（\cref{prop:g7:oxygen-in-the-water:map}）。货币供应：
\omterm{def:g6:exploring-our-environment:conditions}{温度}决定水能容纳多少\omterm{def:g7:what-is-respiration:respiration}{氧气}——低温
提高溶氧上限，温暖则使其降低
（\cref{prop:g7:oxygen-in-the-water:drains}）。最大的
开支者：有死亡物质时，\omterm{def:g6:decomposers-and-soil:decomposers}{分解者}集中的
\omterm{def:g7:what-is-respiration:respiration}{呼吸作用}耗氧最多——奶制品厂排水管
的教训（\cref{ex:g7:oxygen-in-the-water:waste}）。
\end{solution}

\begin{solution}{pb:g7:oxygen-in-the-water:1}
\textbf{1.} K1 水温低——溶氧上限高——且被
堰坝搅动——水面补氧全速进行：读数接近上限。
K2 水温高——溶氧上限降低——且未经搅动：
其供氧来源所能维持的只是一个较低的读数。

\textbf{2.} K1 中要求高的\omterm{def:g5:biodiversity:species}{物种}证明水质富氧、稳定；
K2 中耐受性强的\omterm{def:g5:biodiversity:species}{物种}——鲤鱼、螺类、带\omterm{def:g7:what-is-respiration:respiration}{呼吸}管的幼虫——
表明该水体偏贫氧，至少在其糟糕的时段如此。

\textbf{3.} \omterm{def:g6:decomposers-and-soil:decomposers}{分解者}的\omterm{def:g7:what-is-respiration:respiration}{呼吸作用}：奶制品厂
富含\omterm{def:g6:plants-are-producers:matter}{有机物}的排水养活了大量\omterm{def:g6:decomposers-and-soil:decomposers}{分解者}，它们
的大量进食使\omterm{def:g7:oxygen-in-the-water:dissolved}{溶解氧}存量远低于仅由\omterm{def:g6:exploring-our-environment:conditions}{温度}
所能解释的水平。

\textbf{4.} 两公里的消耗已耗尽了
废物——盛宴已结束，耗氧者也已减少——
而多石的浅滩又把\omterm{def:g7:what-is-respiration:respiration}{氧气}搅回水中：
两个供应者都已恢复，耗氧已停止，
普查证实水体已经恢复。

\textbf{5.} 日间波动：黎明低谷为 \qty{4}{mg/L}，
午后高峰为 \qty{9}{mg/L}。两类“交易者”：白天的光合作用
（白天增氧）与全天候的\omterm{def:g7:what-is-respiration:respiration}{呼吸作用}
（整夜单独耗氧）。

\textbf{6.} 冬季低温抬高了溶氧上限，因此水能容纳更多\omterm{def:g7:what-is-respiration:respiration}{氧气}；
低温也减缓了每一个耗氧者——\omterm{prop:g3:sorting-living-things:groups}{鱼类}、\omterm{def:g1:plants-around-us:plant}{植物}与泥中的\omterm{def:g6:decomposers-and-soil:decomposers}{分解者}——的代谢，
所以耗氧量减少。水质也更稳定，
因为\omterm{def:g5:food-webs:roles}{生产者}的日间波动较小。

\textbf{7.} K3 最先（已耗尽，且热量加速其
\omterm{def:g6:decomposers-and-soil:decomposers}{分解者}的活动）；K2 次之（溶氧上限低，水流静止，
黎明时出现低谷）；K4 第三（正在恢复，但位于问题河段的下游）；
K1 最后（水冷且翻搅——堰坝不断补充\omterm{def:g7:what-is-respiration:respiration}{氧气}，
无论天气如何）。

\textbf{8.} 日出前是每日最低点：读数记录的是
河流生命必须承受的最坏情况。若安排在下午，
就会把全年数据记在它最好的
时刻之下。

\textbf{9.} 探头数据在数周内即可恢复——耗氧随\omterm{def:g6:decomposers-and-soil:decomposers}{分解者}食物耗尽而停止。
\omterm{def:g6:exploring-our-environment:components}{生物}普查恢复较慢，且有先后次序：
先是耐受性强的类群减少，
随后是漂移和飞行而来的定居者
（\cref{ch:g6:how-plants-spread} 的逻辑，\omterm{def:g1:animals-around-us:animal}{动物}版本）
重新\omterm{prop:g6:how-plants-spread:travel}{定殖}要求苛刻的\omterm{def:g5:biodiversity:species}{物种}——蜉蝣幼虫先从上游来，
鳟鱼最后。

\textbf{10.} 树荫也使河段保持凉爽——而低温正是溶氧上限。
砍伐会增加白天的光合产氧，
却会降低溶氧上限，并加速升温水体中的耗氧过程；
在缓流河段，损失很可能超过收益。
答案：保留树木，改为搅动水体。

\textbf{11.} 普查：蜉蝣幼虫无法伪装出良好的季节——
它们经历了其中的每一个黎明。午后的平庸读数只是一个瞬间，
也许是异常的一天；
幼虫才是这一季可信的记录。

\textbf{12.} “让河水保持低温、充分搅动且未受污染，
\omterm{def:g7:what-is-respiration:respiration}{氧气}——以及所有依赖它\omterm{def:g7:what-is-respiration:respiration}{呼吸}的生命——
自会得到照顾。”
\end{solution}
